\chapter{Introducción}
\label{cap:Introduccion}

La codificación clínica representa un proceso crítico en la gestión sanitaria moderna, actuando como interfaz entre la práctica clínica, la administración hospitalaria y la investigación médica mediante el análisis estadístico. 

Este trabajo se centra en el desarrollo de soluciones de codificación automática mediante modelos BERT multilabel, para CIE-10, abordando tres aspectos fundamentales:
\begin{itemize}
\item La creciente complejidad de los sistemas de clasificación médica.
\item Los desafíos operativos en entornos clínicos reales.
\item Las oportunidades tecnológicas emergentes en Procesamiento del Lenguaje Natural (PLN).
\end{itemize}

\section{Motivación}

La motivación de este trabajo surge de varios factores clave identificados:
\begin{itemize}
\item Es necesario disponer de conocimientos medicos para la correcta clasificación, haciendo necesario tener personal especializado.
\item La normativa CIE-10 contiene 68,069 códigos posibles, haciendo que la tarea de clasificar tenga mucha complejidad, y aumentando el número de errores en el proceso de clasificación manual.
\item Desde el 1 de enero de 2016 la CIE-10-ES ha sido implantada en España como clasificación de referencia para la codificación clínica, como indica el Real Decreto 69/2015, y su obligatoriedad se esta implantando de manera gradual
\item El tiempo medio de codificación de un alta hospitalaria es de 14 minutos, 10 minutos más en comparación con la CIE9, tal y como se indica en el informe de implantación de la CIE 10\cite{sanidad2016cie10}
\end{itemize}

La solución propuesta combina técnicas de Deep Learning con conocimiento clínico estructurado, permitiendo:
\begin{itemize}
\item La reducción del tiempos de codificación, asi como la automatización del proceso, requiriendo unicamente la posterior revisión humana en el caso de que la certeza sea baja
\item La consistencia en la aplicación de normas CIE-10 mediante aprendizaje automático
\item La detección automática de inconsistencias en narrativas clínicas
\end{itemize}


\section{Contexto disciplinar y tecnológico}

El análisis de trabajos previos revela tres generaciones tecnológicas:

\textbf{Sistemas basados en reglas} (1990-2010): Utilizan ontologías médicas como SNOMED-CT con lógica booleana. Estos están muy limitados por su alta rigidez ante la variabilidad lingüística~\cite{aronson2011}.

\textbf{Modelos estadísticos} (2010-2018): Aplicación de SVM y Random Forests sobre características léxico-sintácticas. Logran un F1=0.79 en CIE-9 pero no consiguen escalar su poder de clasificación a la CIE-10~\cite{liu2018}.

\textbf{Deep Learning} (2018-actual): Modelos secuencia-a-secuencia con atención contextual. El sistema CodiEsp~\cite{miranda2020} alcanza un F1=0.85 en español usando el modelo neuronal BiLSTM.

\subsection{Avances en transformers médicos}
Los modelos BERT adaptados al dominio clínico, como ClinicalBERT~\cite{alsentzer2019} y BioBERT~\cite{lee2020}, demuestran una gran superioridad en:

\begin{itemize}
\item Capturar relaciones contextuales en narrativas médicas.
\item Manejar ambigüedades terminológicas (ej; 'EPOC' como entidad vs. comorbilidad \\TODO explicar\\).
\item Transfer-learning entre especialidades médicas.
\end{itemize}

Nuestro trabajo extiende estos enfoques mediante:
\begin{enumerate}
\item Adaptación multilingüe para textos clínicos en español
\item Mecanismo de atención jerárquica para códigos padre/hijo en CIE-10.
\item Función de perdida (Loss function) adaptativa.
\item Análisis de pesos en la distribución de los códigos
\end{enumerate}

\subsection{Retos pendientes}
La literatura identifica tres desafíos no resueltos:
\begin{itemize}
\item \textbf{Multilabel y Multiclass extremo}: La CIE-10 contiene 68,069 códigos posibles, con un promedio de 6,1 códigos de diagnostico y 2,7 de procedimientos, esto hace muy complejo el entrenamiento, requiriendo un conjunto de datos muy amplio para poder clasificar de forma consistente cada diagnostico.
\item \textbf{Distribución de los códigos en la codificación}: La distribución de los códigos presenta una gran asimetria, tanto en su jerarquia, como la frecuencia en los diagnosticos médicos, habiendo clases y familias de códigos cuya presencia es mayoritaria, y otras cuya presencia es muy baja, como es el caso de enfermedades raras con una incidencia mínima en la población.
\item \textbf{Evaluación clínica}: Métricas tradicionales (F1, precisión) no capturan las consecuencias clínicas de errores, o el grado de relación entre el código que se obtiene en el clasificador, y el código correcto, por lo que es necesario también tener en cuenta esta limitación e implementar funciones de perdida especificas para este caso. 
\item \textbf{Explicabilidad}: Todo software ligado a la medicina tiene el requisito legal de seguir el Reglamento UE 2017/745. \\TODO desarrollar\
\end{itemize}










